\documentclass[11pt]{article}
\usepackage{fontspec}
\usepackage{polyglossia}
\setmainlanguage{russian}
\setotherlanguage{english}
\setmainfont{Times New Roman}
\newfontfamily\cyrillicfont{Times New Roman}
\newfontfamily\cyrillicfontsf{Helvetica}
\newfontfamily\cyrillicfonttt{Menlo}
\usepackage{amsmath,amssymb}
\usepackage{booktabs}
\usepackage[margin=1in]{geometry}
\usepackage{enumitem}
\usepackage{hyperref}

\title{Технический разбор PhAIL: всё, что нужно знать чтобы защитить статью}
\author{}
\date{}

\begin{document}
\maketitle

Этот документ~--- спутник к статье. Идея: пройти пошагово от "зачем всё это" до "что отвечать рецензенту", объясняя каждый технический термин в момент его появления. Если рецензент спросит "а что такое X?"~--- ответ должен быть здесь.

\tableofcontents
\bigskip

%==============================================================
\section{Зачем всё это: проблема с бинарной долей успехов}
%==============================================================

Стандартная практика в VLA-статьях: "наша модель завершает 70\% эпизодов за $\tau$ секунд". Это прокси надёжности~--- говорит, насколько модель устойчива. Но:

\begin{enumerate}[leftmargin=2em, itemsep=2pt]
\item \textbf{Не учитывает скорость работы.} 90\% за 30~с и 90\% за 300~с~--- абсолютно разные политики, но бинарная доля одинакова.
\item \textbf{Зависит от выбора $\tau$.} На разных $\tau$ ранжирование моделей может меняться (наша Таблица~3 показывает: $F(30s)$ и $F(60s)$ дают разные top-1).
\item \textbf{Нечувствителен к форме распределения.} Политика "или быстро или никогда" и политика "всегда среднеспешная" могут иметь одинаковый $F(\tau)$.
\end{enumerate}

Решение: измерять \emph{распределение} времени до успеха $T$ целиком, а не только его значение в одной точке.

%==============================================================
\section{Случайная величина $T$ и её CDF}
%==============================================================

\paragraph{Определение $T$.} $T$~--- время от начала операции до её успешного завершения. "Операция"~--- это перемещение одного предмета из исходного контейнера в целевой. Один rollout (запуск модели на эпизоде) содержит несколько операций (батарейки~--- 10+, ножницы~--- 1).

\paragraph{Три класса исходов операции.}
\begin{enumerate}[leftmargin=2em, itemsep=2pt]
\item \textbf{Успех}: модель положила предмет в целевой контейнер. $T$~--- время от старта до этого момента, событие $E=1$.
\item \textbf{Правоцензурированный} (right-censored): эпизод закончился по таймауту (30~с/предмет, или оператор остановил), модель ещё работала, но предмет не положен. Знаем $T > T_{\text{tail}}$, но не само $T$. Событие $E=0$.
\item \textbf{Жёсткий провал} (hard failure): предмет уронен за пределы рабочего пространства, или сработала аварийная остановка. Операция \emph{не может} завершиться никогда. Записываем как $T=\infty$, $E=1$ ("событие-призрак", ghost event).
\end{enumerate}

\paragraph{CDF.} $F(t) = P(T \le t)$~--- доля операций, завершённых к моменту $t$. Из определения: $F(0)=0$, $F$ монотонно растёт, $F(\infty) \le 1$ (асимптота ниже~1~$=$~доля жёстких провалов).

\paragraph{Почему CDF, а не функция выживания $S(t)=1-F(t)$.} Эквивалентно с математической точки зрения. CDF интуитивнее для не-survival-аудитории: "доля завершённых" читается напрямую с графика. Функция выживания~--- стандарт в медицине, но мы пишем для ML-сообщества.

%==============================================================
\section{Kaplan-Meier: как оценивать CDF при наличии цензурирования}
%==============================================================

\paragraph{Почему не просто эмпирическая CDF.} Эмпирическая CDF $\hat F(t) = \frac{1}{n}\sum_i \mathbb{1}[T_i \le t]$ требует знания $T$ для всех наблюдений. У нас часть~--- правоцензурированные ($T > T_{\text{tail}}$, точное значение неизвестно). Их нельзя ни выбросить (систематический сдвиг к коротким временам), ни считать как $T = T_{\text{tail}}$ (заниженная оценка).

\paragraph{Оценка Kaplan-Meier (1958).} Оценивает $S(t) = 1 - F(t)$ как произведение условных вероятностей "выжить через момент $t_i$":
\[
\hat S(t) = \prod_{t_i \le t} \left(1 - \frac{d_i}{n_i}\right)
\]
где $t_i$~--- моменты событий (успехов), $d_i$~--- число успехов в момент $t_i$, $n_i$~--- размер выборки-в-риске (тех, кто ещё не "вышел" до $t_i$). Цензурированные наблюдения остаются в $n_i$ до момента цензурирования, потом выпадают~--- без подделки события.

\paragraph{Дисперсия по Greenwood.} Оценка дисперсии Kaplan-Meier-оценки $\hat S(t)$:
\[
\widehat{\text{Var}}(\hat S(t)) = \hat S(t)^2 \sum_{t_i \le t} \frac{d_i}{n_i(n_i - d_i)}
\]
\emph{Откуда формула.} Поскольку $\hat S(t) = \prod (1 - d_i/n_i)$, логарифм превращается в сумму:
\[
\log \hat S(t) = \sum_{t_i \le t} \log\!\left(1 - \frac{d_i}{n_i}\right).
\]
Каждое слагаемое~--- условная вероятность "выжить через момент $t_i$", биномиальная по природе с дисперсией $d_i/(n_i(n_i - d_i))$. Сумма (под предположением независимости наблюдений) даёт $\text{Var}(\log \hat S)$. Возвращаемся к $\hat S$ через delta-метод~--- получаем формулу выше.

\emph{Что она даёт.} Для каждого момента $t$~--- стандартную ошибку $\hat S(t)$. Из неё строится поточечный 95\% доверительный интервал на $S(t)$ (или эквивалентно $F(t) = 1 - S(t)$).

\paragraph{Log-log преобразование.} Если строить наивный 95\% CI как $\hat S(t) \pm 1.96 \cdot \text{SE}$~--- при низких $n$ интервал может выходить за $[0,1]$. Например, $\hat S = 0.05$, $\text{SE} = 0.04$~$\Rightarrow$ нижняя граница $-0.03$. Это бессмысленно для вероятности.

\emph{Решение.} Преобразование $g(S) = \log(-\log S)$ отображает $S \in (0,1)$ в $\mathbb{R}$. Строим симметричный гауссовский CI в трансформированной шкале, потом возвращаемся обратно:
\[
\text{CI}_{95\%}(S(t)) = \hat S(t)^{\exp\big(\pm 1.96 \cdot \text{SE}(\log(-\log \hat S(t)))\big)}
\]
Гарантированно остаётся в $(0,1)$ и асимптотически корректный. Полученный интервал \emph{асимметричен}: при $\hat S$ близком к нулю или единице он шире со стороны, удалённой от границы~--- ровно потому, что вблизи границы вероятность плотно "упирается".

\paragraph{Где Greenwood в нашей статье~--- и где НЕ.} Прямой ответ: \textbf{в главных таблицах и фигурах мы Greenwood не используем}. Все 95\% CI в Таблицах~2--4 и Рисунках~1--5 построены через бутстрап по эпизодам.

\emph{Почему не Greenwood.} Greenwood предполагает независимость наблюдений (операций). У нас операции внутри эпизода коррелированы~--- общая сцена, освещение, состояние политики. Greenwood в этом случае дал бы искусственно узкие интервалы~--- ту же проблему, что и стандартный бутстрап по операциям. Бутстрап по эпизодам корректно учитывает корреляцию.

\emph{Где упоминается.} В \S2.1 статьи: \emph{"pointwise log-log intervals using the Greenwood variance estimator are equivalent and described in Kalbfleisch--Prentice"}~--- то есть мы признаём Greenwood как канонический метод survival analysis, но используем бутстрап как более подходящий для нашего случая. Это покрывает рецензента, который скажет: "почему не стандартный Greenwood?"~--- ответ: коррелированные наблюдения, см.\ \S2.2.

\emph{Когда Greenwood и бутстрап эквивалентны.} Если бы у нас был один эпизод $=$ одна операция (как в медицинском survival, где пациент $=$ одно событие), Greenwood и бутстрап по эпизодам дали бы практически одинаковые CI. Но у нас один эпизод $=$ много операций (батарейки~--- 10+), и эта внутри-эпизодная корреляция критична.

\paragraph{Призраки ($T=\infty$).} Простое расширение: жёсткий провал добавляет дискретный скачок CDF в правом крае. Эквивалентно: для интегральных метрик (RMST) подставляем $T_{\text{ghost}} = \tau$ (потолок интегрирования). Жёсткий провал увеличивает RMST ровно на $\tau$ минус то время, которое политика потратила бы на эту операцию при успехе.

%==============================================================
\section{Скаляры из CDF: что мы публикуем}
%==============================================================

\paragraph{Все стандартные скаляры~--- проекции $F$.}
\[
\begin{aligned}
F(\tau) &= \text{доля успехов к моменту}\ \tau \\
\inf\{t: F(t) \ge 0.5\} &= \text{медианное время до успеха} \\
\int_0^\tau (1 - F(t))\,dt &= \text{RMST}(\tau)\ \text{(restricted mean survival time)} \\
3600 / E[T \mid T < \tau] &\propto 1/\text{RMST} = \text{UPH (units per hour)} \\
\text{RMST}(\tau) / (1 - F(\tau)) &= \text{MTBF/A}\ \text{(mean time between failures/assists)}
\end{aligned}
\]

\paragraph{HRT (Human-Relative Throughput) -- наш главный скаляр.}
\[
\text{HRT}(m, o) = \frac{\text{RMST}_{\text{teleop}, o}(\tau)}{\text{RMST}_{m, o}(\tau)}, \qquad \text{HRT}(m) = \frac{1}{|O|}\sum_{o \in O} \text{HRT}(m, o)
\]
Поячеечное отношение модели к телеоп-референсу, макро-усреднённое по объектам. Свойства:
\begin{itemize}[leftmargin=2em, itemsep=1pt]
\item \textbf{Безразмерность}: можно сравнивать разнородные задачи.
\item \textbf{Поячеечная нормализация}: трудный объект замедляет телеоп и модель пропорционально~--- сложность объекта сокращается.
\item \textbf{$T=\infty$ увеличивает знаменатель}: жёсткие провалы автоматически штрафуют HRT.
\item \textbf{Макро vs микро}: макро-усреднение по объектам, чтобы один тяжёлый объект (мало эпизодов) не доминировал.
\end{itemize}

\paragraph{Почему $\tau = 240$~с в главной таблице.} 30~с/предмет $\times$ 8 предметов в самой длинной задаче $=$ 240~с~--- естественный потолок эпизода в нашем протоколе. Цензурируем $T$ при $\tau$~--- это "сколько модель потратила бы, если бы работала вечно (но мы не наблюдали)".

%==============================================================
\section{Доверительные интервалы: episode-clustered bootstrap}
%==============================================================

\paragraph{Зачем бутстрап.} Хотим доверительный интервал на любой скаляр (RMST, HRT, $F(\tau)$). Аналитические формулы (Greenwood и т.п.)~--- только для отдельных метрик и при предположениях независимости. Бутстрап~--- универсальный непараметрический метод (Efron, 1979).

\paragraph{Идея бутстрапа.} Многократно ($B$ раз) ресэмплируем из наших данных с возвращением, считаем метрику на каждом ресэмпле, берём 2.5-й и 97.5-й перцентили~--- это и есть 95\% доверительный интервал.

\paragraph{Почему \emph{по эпизодам кластеризованный}.} Стандартный бутстрап~--- ресэмпл по \emph{наблюдениям} (т.е.\ по операциям). Но операции внутри одного эпизода \emph{коррелированы}~--- они делят:
\begin{itemize}[leftmargin=2em, itemsep=1pt]
\item Состояние политики (после первой операции изменился контекст).
\item Визуальные условия (та же сцена, освещение, камера).
\item Стартовое состояние робота.
\end{itemize}
Если ресэмплировать по операциям, получаем \emph{искусственно узкие} доверительные интервалы~--- они занижают неопределённость, потому что не учитывают эту корреляцию. Эпизод~--- естественный независимый блок (разные стартовые условия, независимые запуски).

\paragraph{Алгоритм.}
\begin{enumerate}[leftmargin=2em, itemsep=1pt]
\item Берём $n$ эпизодов в группе.
\item Ресэмплируем $n$ эпизодов с возвращением~$\Rightarrow$ бутстрап-выборка.
\item Считаем метрику (RMST, $F(\tau)$, HRT) на этой выборке.
\item Повторяем $B = 1000$--$2000$ раз.
\item Перцентильный интервал: $[\text{quantile}_{2.5},\ \text{quantile}_{97.5}]$.
\end{enumerate}

\paragraph{Парный бутстрап для сравнения двух политик.} При сравнении модели $A$ и модели $B$ (или модели и телеопа) на одних объектах: ресэмплируем \emph{пары} индексов, чтобы сохранить корреляцию между моделями. "Парные" в смысле "один и тот же эпизод-объект в обеих группах", что снижает дисперсию разностей.

%==============================================================
\section{От описания к ранжированию: зачем нужны отдельные тесты}
%==============================================================

\paragraph{Описание (что есть).} "RMST модели~$=$~80~с, 95\%~CI [70, 92]"~--- описательная статистика. Не отвечает на вопрос "отличается ли это от другой модели?".

\paragraph{Наивная ошибка~\#1: "сравнить точки".} "RMST $A=80$, RMST $B=85$, значит $A$ лучше"~--- игнорирует доверительный интервал. Если интервалы пересекаются, разница может быть шумом.

\paragraph{Наивная ошибка~\#2: "пересекаются ли интервалы".} "CI $A$~$=$~[70, 92], CI $B$~$=$~[75, 95], они пересекаются $\Rightarrow$ нет разницы". \emph{Неверно}: пересечение маргинальных интервалов не эквивалентно отсутствию разницы. Правильно сравнивать интервал на \emph{разности} $A - B$, а не два отдельных интервала.

\paragraph{Что мы хотим.} Формальный тест: "какова вероятность увидеть наблюдаемую разницу или больше при $H_0: F_A = F_B$?". Это $p$-значение. Если $p < 0.05$~--- отвергаем $H_0$.

\paragraph{Двое таких тестов в нашей статье.}
\begin{enumerate}[leftmargin=2em, itemsep=1pt]
\item \textbf{Logrank} (раздел~\ref{sec:logrank})~--- классика survival analysis, контрольная проверка.
\item \textbf{Kolmogorov-Smirnov} (раздел~\ref{sec:ks})~--- наш основной тест, более экономный по выборке.
\end{enumerate}

%==============================================================
\section{Logrank: классический тест на равенство CDF}\label{sec:logrank}
%==============================================================

Тест на сравнение двух кривых выживаемости. Спрашивает: $H_0: F_a = F_b$.

\paragraph{Механика.}
\begin{enumerate}[leftmargin=2em, itemsep=1pt]
\item Соберите все события (успехи) из обеих групп и упорядочьте по времени: $t_1 < t_2 < \dots < t_K$.

\item В каждый момент события $t_i$ посмотрите на \emph{выборку-в-риске}~--- операции, которые ещё не завершились к этому моменту в обеих группах. Пусть $n_{a,i}, n_{b,i}$~--- размер выборки-в-риске в группах, $d_i$~--- общее число событий в момент $t_i$.

\item Под $H_0$ ожидаемое число событий из группы $a$ пропорционально доле в выборке-в-риске:
\[
E_{a,i} = d_i \cdot \frac{n_{a,i}}{n_{a,i} + n_{b,i}}
\]

\item Сравните наблюдаемое с ожидаемым:
\[
Z = \frac{\sum_i (O_{a,i} - E_{a,i})}{\sqrt{\sum_i V_i}} \sim \mathcal{N}(0, 1)
\]
где $V_i$~--- гипергеометрическая дисперсия в каждом моменте.

\item $Z^2 \sim \chi^2_1$ под $H_0$~--- это и есть logrank-статистика.
\end{enumerate}

\paragraph{Интуиция.} В каждый момент события спрашиваем: "если бы группы были одинаковыми, какая доля сегодняшних успехов должна быть из группы $a$?". Складываем расхождения. Систематический сдвиг $\Rightarrow$ статистика растёт.

\paragraph{Свойства.}
\begin{enumerate}[leftmargin=2em, itemsep=1pt]
\item Использует \emph{всю кривую}, а не одно $F(\tau)$.
\item Нативно обрабатывает цензурирование.
\item Стандарт в медицинской survival analysis с 1966~г.\ (Mantel).
\end{enumerate}

\paragraph{Слабое место.} Logrank \emph{предполагает пропорциональные риски}~--- что отношение hazard-функций двух групп постоянно по времени. Если кривые \emph{пересекаются} (crossover), logrank теряет мощность. На наших данных пересечений нет, но это известное ограничение, и KS этого не требует.

%==============================================================
\section{KS: Kolmogorov-Smirnov, наш основной тест}\label{sec:ks}
%==============================================================

\paragraph{Статистика.}
\[
D = \sup_t \big|F_a(t) - F_b(t)\big|
\]
максимальное вертикальное расстояние между двумя CDF.

\paragraph{Геометрия через P-P plot.} Если параметрически отложить $\big(F_a(t), F_b(t)\big)$ при $t \in [0, \tau]$~--- получим P-P-кривую. Тогда:
\begin{itemize}[leftmargin=2em, itemsep=1pt]
\item \textbf{AUC} (площадь под кривой) $=$ $P(T_a < T_b)$~--- попарное стохастическое доминирование.
\item \textbf{KS} $=$ максимальное вертикальное расстояние от кривой до диагонали.
\end{itemize}
Обе метрики читаются с одной картинки.

\paragraph{Что говорит и не говорит.} KS отвечает "отличаются?"~--- да или нет с $p$-значением. \emph{Не говорит, какая модель быстрее}. Направление берём из RMST или знака $F_a(t^*) - F_b(t^*)$ в точке максимума. В статье публикуем оба: KS ("отличаются?") $+$ RMST ("кто быстрее?").

\paragraph{Почему НЕ используем стандартное $p$-значение KS.} Учебная асимптотика KS предполагает независимость наблюдений. У нас операции внутри эпизода коррелированы. Поэтому:
\[
p\text{-значение KS} = \text{бутстрап с кластеризацией по эпизодам}
\]
\begin{enumerate}[leftmargin=2em, itemsep=1pt]
\item Объединяем эпизоды обеих групп в общий пул под $H_0$.
\item Многократно (200--2000) ресэмплируем целые эпизоды с возвращением, разбиваем случайно на две группы того же размера, считаем $D^*$.
\item $p$-значение $=$ доля бутстрапа, где $D^* \ge D_{\text{наблюдённое}}$.
\end{enumerate}

\paragraph{Почему KS экономный по выборке.} KS чувствителен к точке, где кривые расходятся максимально. RMST $= \int_0^\tau (1-F(t))\,dt$~--- интегратор, размазывает информацию о расхождении по всему $[0,\tau]$. На пограничных парах KS детектирует разницу при $N=20$, RMST и $F(60s)$~--- никогда при $N \le 30$. Информация живёт в \emph{том, где} кривые разделяются; интегратор её сглаживает, sup-норма её ловит.

%==============================================================
\section{Поправка на множественные сравнения: Bonferroni}
%==============================================================

\paragraph{Какие именно 6 гипотез.} Для четырёх моделей $\binom{4}{2} = 6$ пар. Для каждой пары $(a, b)$:
\[
H_0^{(a,b)}: F_a = F_b \quad\text{против}\quad H_1^{(a,b)}: F_a \neq F_b.
\]
Конкретно:
\begin{enumerate}[leftmargin=2em, itemsep=0pt]
\item OpenPI $=$ GR00T
\item OpenPI $=$ ACT
\item OpenPI $=$ SmolVLA
\item GR00T $=$ ACT
\item GR00T $=$ SmolVLA
\item ACT $=$ SmolVLA
\end{enumerate}

\paragraph{Зачем коррекция: семейная ошибка первого рода (FWER).} При $K$ независимых тестов на уровне $\alpha = 0.05$:
\[
\text{FWER} = P(\text{хотя бы один ложно-положительный}) = 1 - (1-\alpha)^K \approx K\alpha
\]
При $K=6$: FWER~$\approx$~0.26. То есть 26\% шанс опубликовать ложно-значимое сравнение, даже если все модели одинаковы.

\paragraph{Bonferroni.} Тестировать каждую гипотезу на уровне $\alpha/K$. Эквивалентно: $p_{\text{Bonf}} = \min(K \cdot p_{\text{raw}},\ 1)$. Гарантирует FWER~$\le \alpha$ при любой структуре зависимости тестов.

\paragraph{Альтернативы и почему не они.}
\begin{itemize}[leftmargin=2em, itemsep=1pt]
\item \textbf{Без коррекции}: 26\% шанс ложного срабатывания~--- нечестно.
\item \textbf{Holm step-down}: тоже контролирует FWER, чуть менее консервативен. На $K=6$ выигрыш мал.
\item \textbf{Benjamini-Hochberg (FDR)}: контролирует ожидаемую долю ложных открытий, не вероятность хотя бы одного. Для исследовательских анализов с тысячами тестов. Не для несущих ranking-утверждений в статье.
\end{itemize}
Наш контекст: 6 подтверждающих ranking-утверждений, каждое публикуется отдельно. Bonferroni~--- правильный выбор.

\paragraph{Что выживает Bonferroni у нас.}
\begin{itemize}[leftmargin=2em, itemsep=1pt]
\item Все 3 пары vs SmolVLA: $p_{\text{Bonf}} < 10^{-14}$~--- железные \texttt{***}.
\item GR00T vs ACT: $p_{\text{raw}}=2.3\!\cdot\!10^{-3} \Rightarrow p_{\text{Bonf}}=1.4\!\cdot\!10^{-2}$~--- выживает, \texttt{*}.
\item OpenPI vs ACT: $p_{\text{raw}}=1.1\!\cdot\!10^{-2} \Rightarrow p_{\text{Bonf}}=6.7\!\cdot\!10^{-2}$~--- \emph{не} выживает, пограничный.
\item OpenPI vs GR00T: $p_{\text{raw}}=0.45$~--- ни до, ни после.
\end{itemize}
Bonferroni "съел" OpenPI vs ACT~--- без коррекции мы бы заявили значимость, после~--- честно говорим маргинально. В этом и смысл коррекции.

%==============================================================
\section{Анализ мощности: откуда берётся 600--1500 парных rollout'ов}
%==============================================================

\paragraph{Мощность теста.} Вероятность отвергнуть $H_0$, когда $H_1$ верна. Конвенция $1-\beta = 0.8$ (80\% мощность).

\paragraph{Тест McNemar (парный бинарный).} Для пары моделей с парными бинарными исходами (успех/неудача) рассматриваются \emph{несогласные пары}~--- эпизоды, где одна модель преуспела, другая нет. Статистика:
\[
\chi^2 = \frac{(b - c)^2}{b + c}
\]
где $b, c$~--- числа несогласных пар в обоих направлениях. Парность учитывает корреляцию через одинаковую сцену и объект~--- чувствительность выше, чем у непарного теста.

\paragraph{Сколько парных rollout'ов нужно?} Для детекции разницы в доле успехов $\Delta = 0.05$ (5~процентных пунктов) при мощности~$=$~0.8, $\alpha=0.05$, доле несогласных пар $0.2$:
\[
N \approx \frac{(z_{\alpha/2}\sqrt{2p_d} + z_\beta\sqrt{p_d - \Delta^2})^2}{\Delta^2} \approx 1280
\]
где $p_d$~--- доля несогласных пар. Диапазон 600--1500~--- при разных предположениях о $p_d$ (от 0.1 до 0.3).

\paragraph{Связь с нашими данными.} Сообщество работает на $N = 10$--$20$ на условие. McNemar говорит: для сравнений топ-VLA (5~п.п.\ разрыв) нужно 600--1500. Разрыв в 30--75$\times$. Это и есть "two orders of magnitude under-budgeted" из абстракта.

\paragraph{Парный vs непарный тест.}
\begin{itemize}[leftmargin=2em, itemsep=1pt]
\item \textbf{Непарный}: две независимые группы выбираются случайно из своих популяций. Дисперсия выше, нужно больше $N$.
\item \textbf{Парный}: две группы матчатся по релевантным факторам (объект, конфигурация, сцена). Снимает межусловную дисперсию, эффективно требует меньше $N$.
\end{itemize}
В PhAIL "парный" $=$ "один и тот же объект, одна и та же конфигурация камеры/контейнера, разные модели в случайной чередующейся ротации". Каждая (модель, объект, конфигурация)-ячейка~--- условие парности.

%==============================================================
\section{Эксперимент с подвыборками: как мы строим графики мощности}
%==============================================================

McNemar даёт аналитический ответ "сколько надо при бинарной метрике". Но для KS аналитики нет, и мы хотим напрямую сравнить четыре метрики ($F(30s)$, $F(60s)$, RMST, KS) на одних данных. Поэтому делаем \emph{эмпирический эксперимент мощности}: симулируем, как часто каждая метрика даёт $p < 0.05$ при разных размерах выборки.

\paragraph{Что такое detection rate.} Это эмпирическая мощность теста: при $N$ эпизодах на (модель, объект)-ячейку~--- какая доля экспериментов отвергнет $H_0: F_a = F_b$? Если detection rate $=0.8$~--- значит в 80\% случаев тест бы поймал разницу. Конвенциональная цель: $\text{detection rate} \ge 0.8$.

\paragraph{Что такое размер подвыборки $N$.} У нас в датасете 30--45 эпизодов на ячейку. "Подвыборка $N$" $=$ "а что было бы, если бы у нас было только $N$?". Тестируем $N \in \{5, 10, 15, 20, 25, 30\}$. $N=5$~--- симулирует типичную VLA-статью; $N=30$~--- ближе к нашему фактическому размеру.

\paragraph{Алгоритм для одной точки на графике.} Зафиксируем пару моделей (например, GR00T vs ACT), метрику (например, KS) и $N=20$.

\begin{enumerate}[leftmargin=2em, itemsep=2pt]
\item \textbf{Внешний цикл~--- 300 trials}. Каждый trial симулирует один "паралельный мир":
   \begin{itemize}[leftmargin=1.5em, itemsep=0pt]
   \item Для каждой (модель, объект)-ячейки берём подвыборку из $N=20$ эпизодов с возвращением. Получаем 4 ячейки $\times$ 2 модели $=$ 8 подвыборок по 20 эпизодов.
   \item Считаем $p$-значение KS на этих подвыборках (см.\ внутренний цикл ниже).
   \end{itemize}
\item \textbf{Внутренний цикл (внутри одного trial)~--- бутстрап в 200 повторов}, чтобы получить $p$-значение KS:
   \begin{itemize}[leftmargin=1.5em, itemsep=0pt]
   \item Считаем наблюдаемую KS-статистику $D_{\text{obs}}$ на подвыборках двух моделей.
   \item 200 раз: объединяем подвыборки в общий пул под $H_0$, ресэмплируем целые эпизоды, разбиваем случайно на две группы того же размера, считаем $D^*$.
   \item $p = $ доля из 200, где $D^* \ge D_{\text{obs}}$.
   \end{itemize}
\item \textbf{Detection rate}~$=$~доля из 300 внешних trials, где $p < 0.05$. Это и есть точка на графике для этой (пары, метрики, $N$).
\end{enumerate}

\paragraph{Что получается.} Для каждой пары моделей (6 пар) и каждой метрики (4 метрики) и каждого $N$ (6 значений) $=$ 144 точки. Складываем в кривые detection rate vs $N$~--- по одной кривой на каждую (пара, метрика).

\paragraph{Как читать график (Рисунок~4 в статье).}
\begin{itemize}[leftmargin=2em, itemsep=1pt]
\item Ось $x$: $N$ (5--30).
\item Ось $y$: detection rate (0--1).
\item Кривая на $y=1$ означает "тест всегда детектирует разницу"~--- эффект огромный.
\item Кривая около $y=0.05$~--- детектирует с шумом (примерно как от случайности).
\item Пересечение $y=0.8$~--- конвенциональная цель 80\% мощности.
\end{itemize}

\paragraph{Что мы видим на трёх режимах.}
\begin{itemize}[leftmargin=2em, itemsep=1pt]
\item \textbf{Насыщенные пары} (любая модель vs SmolVLA): все четыре метрики выходят на $y=1$ уже при $N \le 10$. Эффект огромный, любой тест поймает.
\item \textbf{Близкие пары} (top-3 модели между собой): только KS поднимается. На GR00T vs ACT KS пересекает 0.8 при $N=20$. Остальные метрики остаются на полу ($<0.30$ при $N=30$).
\item \textbf{Истинная ничья} (OpenPI vs GR00T): даже KS не достигает 0.8 при $N=30$ (поднимается до 0.61). Остальные метрики около 0.05~--- неотличимо от случайности.
\end{itemize}

\paragraph{Параметры эксперимента и почему именно такие.}
\begin{itemize}[leftmargin=2em, itemsep=1pt]
\item \textbf{300 внешних trials}~--- даёт Monte-Carlo стандартную ошибку $<3\%$ на detection rate (в худшем случае при $\text{rate}=0.5$).
\item \textbf{200 внутренних повторов}~--- разрешение $p$-значения до $0.005$. При большем числе бутстрап-повторов $p$-значения становились бы точнее, но за рамками 0.05-границы это не меняет detection rate.
\item \textbf{$\tau = 120$\,с} для эксперимента (вместо 240\,с в главной таблице)~--- сужает диапазон интегрирования RMST до области, где модели реально работают, чтобы дать RMST максимально честный шанс. Если бы RMST провалила тест при $\tau=120$\,с, не помогло бы и $\tau=240$\,с.
\end{itemize}

\paragraph{Связь с реальным $N$.} Наши per-cell $N \approx 30$--$45$. Подвыборка $N=30$~--- близко к полному размеру (поэтому KS detection rate при $N=30$~--- хорошая аппроксимация настоящего $p$-значения KS на полных данных). Подвыборка $N=5$~--- симулирует типичную VLA-статью. Разница между двумя крайностями~--- наша headline-картинка о том, как сильно $N$ влияет на разрешающую способность.

%==============================================================
\section{Logrank vs KS на наших данных: согласуются ли?}
%==============================================================

Оба теста на одну и ту же $H_0$, разные статистики. Если \emph{не согласуются}~--- сигнал проблемы. Если согласуются~--- ranking-утверждение надёжно.

\begin{center}
\begin{tabular}{lrrl}
\toprule
Пара & logrank $p_{\text{raw}}$ & KS det.\ rate ($N{=}30$) & Согласие \\
\midrule
OpenPI vs SmolVLA & $<10^{-15}$ & $1.00$ & оба ***  \\
GR00T  vs SmolVLA & $<10^{-15}$ & $1.00$ & оба ***  \\
ACT    vs SmolVLA & $3.3\!\cdot\!10^{-16}$ & $1.00$ & оба ***  \\
GR00T  vs ACT     & $2.3\!\cdot\!10^{-3}$  & $0.95$ & оба значимо \\
OpenPI vs ACT     & $1.1\!\cdot\!10^{-2}$  & $0.75$ & оба пограничные \\
OpenPI vs GR00T   & $0.45$ & $0.61$ & оба не значимо \\
\bottomrule
\end{tabular}
\end{center}

\textbf{Согласие на всех 6 парах.} Качественный рейтинг совпадает. KS более чувствителен на близких парах (отсюда выбор как основной), logrank как контрольная проверка (стандарт литературы).

%==============================================================
\section{Полный workflow ranking-утверждений}
%==============================================================

\begin{enumerate}[leftmargin=2em, itemsep=2pt]
\item \textbf{Объединённая CDF на каждую модель} (по 4 объектам) с бутстрапом по эпизодам для доверительных интервалов.
\item \textbf{Скаляр для отчёта}: RMST или HRT $\pm$ 95\% CI. Не используется для решения "отличаются ли модели".
\item \textbf{Тест для ранжирования}: попарный KS на макро-CDF $+$ бутстрап по эпизодам, 6 тестов, поправка Bonferroni.
\item \textbf{Контрольная проверка}: попарный logrank, тоже Bonferroni, 6 тестов. Сравниваем согласие.
\item \textbf{Направление}: знак $\Delta\text{RMST}$~--- какая модель быстрее.
\end{enumerate}

KS отвечает на "отличаются?", RMST/HRT~--- на "кто быстрее", Bonferroni~--- на "сколько ложных срабатываний мы публикуем". Это и есть "two-jobs framing" статьи: скаляр для отчёта, тест для ранжирования.

%==============================================================
\section{Возможные вопросы рецензентов}
%==============================================================

\paragraph{"Почему не использовали AUC как главный скаляр? Он же даёт top-1 для ACT."} AUC~$=$~$P(T_{\text{model}} < T_{\text{teleop}})$. На наших данных телеоп~$\sim$~7$\times$ быстрее, поэтому AUC доминируется коротковременной областью~--- модель с пиком "быстро-когда-получается" (ACT) выигрывает там, даже если её steady-state хуже. Это свойство референса, не самого AUC. RMST интегрирует по всему $[0,\tau]$ равномерно. Оба корректны; мы выбрали HRT (через RMST) потому что (а) совмещает надёжность и скорость через $T=\infty$, (б) UPH-эквивалент, ближе к operator-практике. \emph{Главное~--- само это разногласие и есть наш эмпирический вклад.}

\paragraph{"$N \approx 35$ маловато даже для вашего теста."} Согласны: на ближайшей паре (OpenPI vs GR00T) даже KS не достигает 80\% мощности при $N=30$. Статья это явно говорит и предлагает $N=600$--$1500$ как требуемый бюджет. Но ключевое утверждение не "мы разрешили все пары"~--- ключевое в том, что (а) GR00T-vs-ACT мы разрешили при $N=20$, (б) бинарные метрики не разрешат этого ни при каком практичном $N$. "The choice of test, not the choice of $N$, determines what is resolvable."

\paragraph{"Почему бутстрап по эпизодам, а не стандартный?"} Операции внутри эпизода коррелированы (общее состояние политики, сцена, освещение). Стандартный бутстрап по операциям даёт искусственно узкие доверительные интервалы. Эпизод~--- естественный независимый блок. Это известная техника (Efron-Tibshirani для кластеризованных данных).

\paragraph{"Почему KS, а не Anderson-Darling или Cramer-von Mises?"} Anderson-Darling сильнее весит хвосты~--- хорошо, когда нас интересуют события на хвостах, у нас не тот случай. Cramer-von Mises использует $L_2$-норму, KS~--- $L_\infty$. Различие на наших данных в пределах шума, но KS~--- стандарт, легче объяснить рецензенту, и его геометрия (максимальное расстояние на P-P-plot) делает его читаемым с рисунка.

\paragraph{"Откуда 600--1500 парных rollout'ов? Покажите расчёт."} Формула McNemar: $N \approx (z_{\alpha/2}\sqrt{2p_d} + z_\beta\sqrt{p_d - \Delta^2})^2/\Delta^2$ при $\Delta=0.05$ (5~п.п.), мощность~0.8, $\alpha=0.05$. Диапазон 600--1500 $=$ при $p_d \in [0.1, 0.3]$. Полный вывод в Приложении~F статьи.

\paragraph{"Почему дефолтные рецепты файн-тюнинга? Это унижает модели."} Намеренно. Оптимизация под каждую модель сместила бы ранжирование и сделала бы статью про конкретные рецепты, а не про методологию. Методология~--- вклад; ранжирование~--- иллюстрация. Статья это явно говорит ("the contribution is methodological").

\paragraph{"Почему 4 объекта? Почему именно эти 4?"} Деревянные ложки (твёрдые, удлинённые), полотенца (деформируемые), ножницы (шарнирные, металлические), батарейки (мелкие, твёрдые)~--- разные физические режимы. 4~--- компромисс между покрытием и бюджетом на аннотацию и телеоп-референс. Бенчмарк спроектирован расти~--- следующие задачи (вставка, мелкая сборка) commited в Future Work.

\paragraph{"Как вы знаете, что аннотации надёжны?"} Двойной поток: высокотировые краудсорс-аннотаторы (Толока) $+$ автоматический пайплайн детекции. Краудсорс~--- золотой стандарт, автоматический оценивается против него на каждом эпизоде. Расхождения публикуются и используются для уточнения автоматического, не наоборот. Оба потока публикуются.

\paragraph{"В Приложении~C GR00T теряет 22~п.п.\ при противоположной стороне камеры. Это разве не главная проблема?"} Это \emph{и есть} вторая методологическая claim: средовые конфаунды могут двигать ранжирование сильнее, чем разрывы между моделями. Поэтому same-fixture рандомизированная ротация со слепым оператором~--- предусловие для всего остального. PhAIL логирует конфигурацию с каждым эпизодом~--- конфаунд виден аналитику, а не растворяется в шуме.

\paragraph{"Почему single-blind, а не double-blind?"} ED-track NeurIPS позволяет single-blind \emph{только} для dataset-centered подач. Наша статья~--- бенчмарк/датасет, поэтому подаём single-blind с открытым именем автора и аффилиацией. Это \emph{требование} call-а, не прихоть.

\paragraph{"Что если рецензент скажет: "just use AUC, it's a single number"?"} Это и есть главный эмпирический вклад: "одно число" зависит от выбора. AUC, RMST, $F(30s)$~--- три principled scalars, три разных top-1. На наших данных это \emph{доказано} (Таблица~3, \S5). Если нет CDF, одно число может ввести в заблуждение; если есть CDF, выбор скаляра~--- осознанный выбор приоритета представления.

\end{document}
